\begin{afterword}
\summaryhead

The post starts from the rule of classical logic that equals may be substituted for equals. Applied to beliefs, the rule gives absurd results. The morning star and the evening star are both Venus; Mary believes the morning star is the god Lucifer; must John, who knows the two are one, conclude that Mary believes the evening star is Lucifer? A simpler version: since ``2 + 2 = 4'' equals TRUE, and so does Fermat's Last Theorem, believing one would mean believing the other.

The author calls this a type error, like subtracting grams from meters. A belief about ``the morning star'' is not the same kind of thing as the morning star, and ``morning star'' is not equal to ``evening star'' even though the planets are one. The cure is to use quotation marks, or a visibly different encoding, wherever beliefs are meant. Mathematical logic and Tarski's sentences about truth train the same habit, and it also separates truth, which compares a belief to reality, from reality itself.

\responsehead

The post's advice, to keep a belief apart from what it is about as a physicist keeps track of units, is sound. The problems are in how the post relates it to earlier work.

The puzzle is Frege's. Frege's 1892 paper used the morning star and the evening star, and the standard account credits Frege with identifying this puzzle about belief reports. The ``simpler version'' is also close to Frege, who held that a sentence stands for its truth value, which is why substitution inside belief reports needed an answer. The post notes only that the problem is ``usually phrased'' this way.

The answer is close to Frege's too. Frege held that after ``believes'' a term stands for a way of conceiving its object, not for the object. The post puts a quoted item in that place. It offers this as its own view (``P'rsnally'', ``in my view''), and the reader is not told how near it is to Frege's.

The AI version, which the post asks the reader to ``imagine'', is on record too. McCarthy wrote in 1979 that ``The problem identified by Frege'' arises ``in artificial intelligence as well as in philosophy and linguistics'', and McCarthy's remedy was the type distinction the post asks for: separate terms for a thing and for the concept of it.

This matters for the post's central claim. It reports that the argument ``is still going on'' and then declares that ``The whole paradox stems from the failure to use quote marks''. But quotation marks were one of the proposals in that argument. McCarthy described that approach as working ``without modal operators or quotation marks and without the restrictions on substituting equals for equals that either device makes necessary.'' The post does not say which proposals it disagrees with, or why quotation succeeds where they fail.

In short: sound advice about keeping beliefs and things apart, applied to a puzzle that is Frege's, with a remedy close to Frege's and McCarthy's, presented as settling an argument whose rival proposals it does not name.
\end{afterword}
